% !PRINTED 79
\weberpage{79}

If the prime $q$ divides $m$ only once and $\mathfrak{A}$ is divisible by $Q$, then $\mathfrak{A}$ is indeed not primary; but the period $\eta$ nevertheless belongs to the group $\mathfrak{A}$.

In this case, however, $\eta$ can be represented by roots of unity of lower degree. This representation may be found in the following way. All numbers of $Q$ are of the form $1 + hm'$, where $h$ runs through the sequence of numbers $0, 1, \dots q - 1$, with the exception of the single value for which

\weq{18}{1 + hm' = kq}

becomes divisible by $q$. If, then, under the assumption that $\mathfrak{A}$ is divisible by $Q$, we set

\[
\mathfrak{A} = Q + Qa_1 + Qa_2 \dots,
\]

it follows that every number $a$ in $\mathfrak{A}$ is of the form:

\weq{19}{a = a'(1 + hm'),}

where $a'$ runs through the sequence of numbers $1, a_1, a_2 \dots$. From this it is further to be concluded that the residues of the numbers $1, a', a'' \dots$ with respect to the modulus $m'$ form a group. If we now take the sum over all values $h = 0, 1, \dots q - 1$ and then remove again the one value determined by (18), it follows:

\[
\eta = \sum^a r^a = \sum^{a'} \sum^h r^{a'(1 + hm')} - \sum^{a'} r^{a'kq},
\]

or, since $\sum^h r^{a'm'h} = 0$:

\[
\eta = -\sum^{a'} r^{a'kq}.
\]

Now $r^q$ is an $m'$th root of unity, and $\eta$ is therefore, apart from the sign, equal to a period formed from $m'$th roots of unity\footnote{Here the memoir of \gesperrt{Fuchs} is to be compared: Ueber die Perioden, welche aus den Wurzeln der Gleichung $\omega^n = 1$ gebildet sind, wenn $n$ eine zusammengesetzte Zahl ist. Crelle's Journal, Bd. 61 (1863).}.

\begin{center}
\S.\,21.

\gesperrt{Cyclotomic Fields with Given Group}.
\end{center}

In the foregoing, then, it has emerged that every cyclotomic period

\[
\eta = \sum^a r^a,
\]

in which $r$ is a primitive $m$th root of unity and $a$ runs through the numbers of a group $\mathfrak{A}$, when $\mathfrak{A}$ is a primary divisor of $\mathfrak{N}$
%CONT>>
